\section{Rekursive Eigenschaft der Kugelvolumen\label{gamma:section:spheres}}
\kopfrechts{Kugelvolumen}

$\pi$ ist definiert als die Konstante, welche 
\eqref{gamma-eq:circumference} erfüllt.

\begin{equation}
\label{gamma-eq:circumference}
    2 \pi r = \text{Umfang eines Kreises mit dem Radius $r$}
\end{equation}

\noindent
Daraus lässt sich die Fläche einer Kreisscheibe herleiten. 
Wenn die Breite eines Rings gegen 0 geht, verhält sich das 
Wachtstum des Flächeninhalts gleich wie das des Umfangs des Rings,
was eine Herleitung der Kreisscheibenfläche per Integral möglich 
macht \eqref{gamma-eq:area}. Für eine grafische Darstellung dieser 
Intuition siehe Abb. \ref{gamma-fig:ring-integration}.

\begin{equation}
\label{gamma-eq:area}
    A(r) = \int_{0}^{r} 2 \pi r \,dr = \pi r^2
\end{equation}

Der Satz des Pythagoras gibt uns eine einfache Beschreibung aller 
Punkte, die Teil des Kreises sind \eqref{gamma-eq:pythagoras}.

\begin{equation}
\label{gamma-eq:pythagoras}
    x^2 + y^2 \leq r^2
\end{equation}

\noindent
Aufgrund der Ungleichung \eqref{gamma-eq:pythagoras} 
lässt sich eine weitere Definition der Kreisfläche herleiten, 
wenn wir sie nach $y$ auflösen \eqref{gamma-eq:pythagoras-y}
und für jedes $x$ den jeweils höchsten $y$-Wert integrieren
\eqref{gamma-eq:area-pythagoras}. Die Intuition dahinter wird
in Abb. \ref{gamma-fig:rect-integration} dargestellt.

\begin{gather}
\label{gamma-eq:pythagoras-y}
    y \leq \sqrt{r^2 - x^2} \\
\label{gamma-eq:area-pythagoras}
    A(r) = 2 \int_{-r}^{r} \sqrt{r^2-x^2} \,dx
\end{gather} 

Das Integral aus \eqref{gamma-eq:area-pythagoras} ist nicht trivial 
zu lösen, deswegen werden wir es vermeiden. Ein Vorteil dieses Ansatzes ist jedoch, 
dass sich die durch den Satz des Pythagoras definierte Metrik sehr 
einfach auf weitere Dimensionen erweitern lässt, zum Beispiel drei Dimensionen 
wie in \eqref{gamma-eq:pythagoras-3d}.

\begin{equation}
\label{gamma-eq:pythagoras-3d}
    \begin{split}
    x^2 + y^2 + z^2 \leq r^2 \\
    z \leq \sqrt{r^2 - x^2 - y^2}
    \end{split}
\end{equation}

Nun können wir das Volumen einer Kugel herleiten, indem wir Halbkreise entsprechender
Grösse entlang $x$ integrieren (siehe Abb. \ref{gamma-fig:semicircle-integration}). 
Dabei integrieren wir entlang der $x$-Achse und innerhalb iterativ entlang der $y$-Achse. 
Eine formelle Begründung für diese iterative angehensweise ist der Satz von Fubini.

\input{papers/gamma/spheres_drawings.tex}

Für ein bestimmtes $x$ sind die möglichen $y$-Werte gemäss Ungleichung 
\eqref{gamma-eq:pythagoras-y} eingeschränkt und für ein bestimmtes $x$ und $y$ wird dann
der maximale Wert für $z$ gemäss Ungleichung \eqref{gamma-eq:pythagoras-3d} genommen. 
Dies ist in \eqref{gamma-eq:sphere-volume} zu sehen. 

Da wir aber bereits die Fläche eines Kreises berechnen können, können wir $A(r)$
zur Berechnung des Kugelvolumens anstatt der Pythagoras-Methode verwenden, wie in 
\eqref{gamma-eq:composed-sphere-volume}.

\begin{equation}
\label{gamma-eq:sphere-volume}
    V(r) = 2 \int_{-r}^{r} \int_{-\sqrt{r^2-x^2}}^{\sqrt{r^2-x^2}} \sqrt{r^2-x^2-y^2} \,dy \,dx
\end{equation}

\begin{equation}
\label{gamma-eq:composed-sphere-volume}
    V(r) = 2 \int_{-r}^{r} A\left(\sqrt{r^2-x^2}\right) \,dx
\end{equation}

Wir haben etwas interessantes getan, nämlich haben wir unsere Berechnung für zwei Dimensionen 
verwendet, um das Volumen in drei Dimensionen zu berechnen. 

Auf diese Weise lässt sich also das Volumen auf beliebige Dimensionen $> 2$ erweitern, und zwar 
mit \eqref{gamma-eq:n-sphere-volume-recursive}.

\begin{equation}
\label{gamma-eq:n-sphere-volume-recursive}
    V_{n+1}(r) = \int_{-r}^{r} V_n\left(\sqrt{r^2-x^2}\right) \,dx
\end{equation}

\noindent
Betrachtet man $A(r) = V_2(r) = 2 \pi r$, so scheint es intuitiv, dass $V_n$ jeweils mit $C_n r^n$
dargestellt werden kann, wobei $C_n$ eine der Dimension zugehörige Konstante ist.

Wir können induktiv beweisen, dass $V_{n+1}$ dieser Form entspricht, sofern
$V_n$ dieser Form entspricht, siehe~\eqref{gamma-eq:const-structure-induction}.

\begin{gather}
\label{gamma-eq:const-structure-induction}
    \text{Nehme an, dass } V_n(r) = C_n r^n \text{. Dann gilt:} \nonumber \\
    \begin{split}
         V_{n+1}(r) &= \int_{-r}^{r} C_n \left(\sqrt{r^2-x^2}\right)^n \,dx  \\ 
        &= \int_{-r}^{r} C_n \left(r^2-x^2\right)^{n/2} \,dx \\
        &= \int_{-1}^{1} C_n \left(r^2 - (ru)^2\right)^{n/2} r \,du \\
        &= \int_{-1}^{1} C_n \left(r^2 (1 - u^2)\right)^{n/2} r \,du \\
        &= \int_{-1}^{1} C_n r^{n+1} (1- u^2)^{n/2} \,du \\
        &= r^{n+1} \underbrace{C_n \int_{-1}^{1} (1- u^2)^{n/2} \,du}_{C_{n+1}} \\
    \end{split}
\end{gather}

Wir sehen, dass $C_n$ ein Faktor von $C_{n+1}$ ist. Die mit dem Volumen assoziierte Konstante 
scheint also eine rekursive Struktur ähnlich der einer Fakultät aufzuweisen. Eine Untersuchung des Terms $I_n = \int_{-1}^1 (1- u^2)^{n/2} \,du$ in \eqref{gamma-eq:n-sphere-constant} ergibt eine weitere rekursive Struktur: $I_n = \frac{n}{n+1} I_{n-2}$.

\begin{equation}
\label{gamma-eq:n-sphere-constant}
    \begin{split}
        I_n &= \int_{-1}^1 (1- u^2)^{n/2} \,du \\
        &= \left[ u (1-u^2)^{n/2} \right]_{-1}^{1} - \int_{-1}^1 u \cdot (-nu(1-u^2)^{n/2 - 1}) \,du \\
        &= (0 - 0) - \int_{-1}^1 -nu^2(1-u^2)^{n/2 - 1} \,du \\
        &= n \int_{-1}^1 u^2 (1 - u^2)^{n/2 - 1} \,du \\
        &= n \int_{-1}^1 (1 - (1 - u^2)) (1-u^2)^{n/2 - 1} \,du \\
        &= n \int_{-1}^1 (1 - u^2)^{n/2 - 1} \,du - n \int_{-1}^1 (1-u^2)^{n/2} \,du \\
        &= n I_{n-2} - n I_n \\
        (n+1) I_n &= n I_{n-2} \\
        I_n &= \frac{n}{n+1} I_{n-2} \\
    \end{split}
\end{equation}

Mit diesem Wissen lässt sich der konstante Faktor des Volumens einer
$n$-Dimensionalen Kugel herleiten als $C_n = C_{n-1} I_{n-1} = C_{n-2} I_{n-1}
I_{n-2} = C_{n-2} \frac{2}{n} I_1 I_0$. Der letzte Schritt ist
durch~\eqref{gamma-eq:n-sphere-constant-2} begründet, funktioniert aber nur mit
geraden $n$.

\begin{equation}
\label{gamma-eq:n-sphere-constant-2}
    \begin{split}
    I_{n-1} I_{n-2} &= (\frac{n-1}{n} \cdot \frac{n-3}{n-2} \cdot \ldots \cdot \frac{3}{4} \cdot I_1) \cdot (\frac{n-2}{n-1} \cdot \frac{n-4}{n-2} \cdot \ldots \cdot \frac{2}{3} \cdot I_0) \\
    &= \frac{n-1}{n} \cdot \frac{n}{n-1} \cdot \ldots \cdot \frac{3}{4} \cdot  \frac{2}{3} \cdot I_1 I_0 \\
    &= \frac{2}{n} \cdot I_1 I_0
    \end{split}
\end{equation}

Von~\eqref{gamma-eq:area} wissen wir, dass $C_2 = \pi$. Eine einfache Rechnung
ergibt, dass $I_0 = 2$ und $I_1 = \frac{\pi}{2}$, d.h. $C_n = C_{n-2} \frac{2 \pi}{n}$. Jetzt können wir für gerade Dimensionen
die rekursive Definition `abrollen' in einen einzigen
Term:~\eqref{gamma-eq:nonrecursive-constant}.

\begin{equation}
\label{gamma-eq:nonrecursive-constant}
    \begin{split}
        C_n &= \pi \cdot \frac{2 \pi}{4} \cdot \frac{2 \pi}{6} \cdot \ldots \cdot \frac{2 \pi}{n} \\
            &= \frac{2 \pi}{2} \cdot \frac{2 \pi}{4} \cdot \frac{2 \pi}{6} \cdot \ldots \cdot \frac{2 \pi}{n} \\
            &= \frac{(2 \pi)^{n/2}}{2^{n/2}} \frac{1}{\left(\frac{n}{2}\right)!} \\
            &= \frac{\pi^{n/2}}{\left(\frac{n}{2}\right)!}
    \end{split}
\end{equation}

Wir können jetzt also wie in \eqref{gamma-eq:n-sphere-volume-even} ein Kugelvolumen bei geraden Dimensionen einfach berechnen.

\begin{equation}
\label{gamma-eq:n-sphere-volume-even}
    V_n(r) = \frac{\pi^{n/2}}{(\frac{n}{2})!}  r^n
\end{equation}

Schön wäre es, das Volumen nicht nur für gerade-Dimensionale Kugeln berechnen können. Dazu müssten wir allerdings die Fakultät auf halbe Zahlen anwenden können.
